\section{Úvod a cieľ projektu}

V súčasnosti zohráva spracovanie prirodzeného jazyka (Natural Language Processing, NLP) dôležitú úlohu pri vývoji inteligentných systémov, ktoré dokážu porozumieť používateľským požiadavkám a poskytovať relevantné odpovede. Jednou z aktuálnych úloh v tejto oblasti je automatické odpovedanie na otázky o tabuľkových dátach formulované v prirodzenom jazyku. Riešenie takýchto úloh si vyžaduje nielen porozumenie textovej otázke, ale aj schopnosť analyzovať štruktúru tabuľky, vyhľadávať potrebné informácie a vykonávať výpočty.

Cieľom tohto projektu je preskúmať možnosti jazykových modelov pri riešení úloh automatického odpovedania na otázky nad tabuľkovými dátami s využitím datasetu DataBench. Vstupom systému je otázka formulovaná v prirodzenom jazyku a príslušná tabuľka. Výstupom je odpoveď vo forme logickej hodnoty (Boolean), kategórie, čísla alebo zoznamu hodnôt. Hlavná pozornosť sa venuje presnosti odpovedí a schopnosti systému správne interpretovať otázky a získavať potrebné informácie z tabuliek.

Hlavnými cieľmi projektu sú:
\begin{enumerate}
    \item analyzovať úlohu DataBench, jej štruktúru a typy otázok a odpovedí,
    \item navrhnúť a implementovať základné riešenie na automatické generovanie odpovedí na otázky týkajúce sa tabuľkových dát,
    \item určiť vhodné metriky hodnotenia a overiť kvalitu základného riešenia na dostupných dátach,
    \item analyzovať výhody a obmedzenia zvoleného prístupu a identifikovať možnosti jeho zlepšenia v ďalšej etape projektu.
\end{enumerate}

\textbf{Výskumná otázka projektu:} Do akej miery dokáže jazykový model efektívne odpovedať na otázky o tabuľkových dátach a aké prístupy môžu zvýšiť presnosť získaných odpovedí?